\documentclass[11pt]{article}
\usepackage{amsmath,amssymb}
\usepackage[margin=1in]{geometry}
\usepackage{microtype}
\usepackage{hyperref}

\newcommand{\dd}{\,\mathrm{d}}
\newcommand{\vv}{\mathbf{v}}
\newcommand{\rd}{\rho_d}
\newcommand{\xhat}{\hat{\xi}}

\title{Where the dry--air mass drift in the moist BF02 benchmark comes from}
\author{Scythe diagnostic note}
\date{\today}

\begin{document}
\maketitle

\begin{abstract}
The moist Bryan--Fritsch (2002) benchmark loses dry--air mass at
$\sim 3\times10^{-3}\,\%$ over $200$\,s, three orders of magnitude more than the
dry case at identical resolution. This note shows, with an offline diagnostic
that uses the model's own spectral operators, that the leak is \emph{not} the
advective form of the continuity equation (its discrete tendency conserves mass
to $\sim10^{-9}\,\%$), but the \textbf{$l_q$ third--derivative smoothing penalty}
in the cubic B--spline least--squares fit, applied to the \emph{nonlinear}
log--density $\xi=\ln(\rd/\rho_0)$ on every timestep. The bare $b_{kDim}<kDim$
projection (with $\mathrm{mubar}=3$ Gauss nodes per cell) is integral--preserving;
it is the smoothing of a nonlinear variable that breaks conservation. Two
independent in--loop tests confirm this: setting $l_q=0$ cuts the loss $120\times$,
and projecting the \emph{linear} density $\rd$ (flux form) conserves mass to
machine precision \emph{with or without} $l_q$. The unbalanced base--state cloud
loading is the trigger only because it drives the circulation that injects the
sub--grid $\xi$ structure the smoothing then removes.
\end{abstract}

\section{Setup}

The dry--air continuity equation is integrated in advective (non--conservative)
form for the log--density perturbation $\xi'=\xi-\xhat$,
\begin{equation}
  \frac{\partial \xi'}{\partial t}
   = -\vv\cdot\nabla\xi' \;-\; w\,\frac{\partial\xhat}{\partial z}
     \;-\; \nabla\cdot\vv ,
  \qquad \rd=\rho_0\,e^{\xi}, \quad \xi=\xi'+\xhat .
  \label{eq:cont}
\end{equation}
Total dry mass is the Gauss--weight quadrature of $\rd$ over the mish points,
$M=\sum_i W_i\,\rd{}_{,i}$, which is the exact integral of the model's cubic
representation. The diagnostic reproduces the benchmark's reported drift
exactly, so $M$ and its drift are correct.

\section{The advective discretization is not the leak}

Continuously, \eqref{eq:cont} is identical to the flux form
$\partial_t\rd=-\nabla\cdot(\rd\vv)$, so $\dd M/\dd t=-\oint\rd\,\vv\cdot\mathbf n=0$
at rigid walls. Discretely one finds, writing $D$ for the spline derivative,
\begin{equation}
  \frac{\dd M}{\dd t}
   = \sum_i W_i\,\rd{}_{,i}\Big(\tfrac{\dd\xi}{\dd t}\Big)_i
   = \underbrace{-\sum_i W_i\big[\nabla\!\cdot(\rd\vv)\big]_i}_{\text{flux form}\;\approx\;0}
     \;+\; \underbrace{\sum_i W_i\,\vv_i\cdot\mathbf E_i}_{\langle\vv,\mathbf E\rangle},
  \qquad
  \mathbf E = D(e^{\xi})-e^{\xi}\,D(\xi).
  \label{eq:defect}
\end{equation}
$\mathbf E$ is the discrete chain--rule defect (zero in the continuum, nonzero
because $e^{\xi}$ is not in the spline space). \emph{This was the original
hypothesis, and it is wrong.} Evaluated from the saved fields with the model's
operators:
\begin{center}
\begin{tabular}{lrr}
\hline
quantity & moist & dry \\
\hline
$\int \dd M/\dd t\big|_{\text{advective RHS}}$ & $5.8\times10^{-10}\,\%$ & $2.2\times10^{-7}\,\%$ \\
$\int \dd M/\dd t\big|_{\text{flux form}}$      & $6.8\times10^{-17}\,\%$ & $6.1\times10^{-18}\,\%$ \\
\textbf{observed drift}                          & $-3.0\times10^{-3}\,\%$ & $-4.3\times10^{-6}\,\%$ \\
\hline
\end{tabular}
\end{center}
The advective RHS conserves mass to $\sim10^{-9}\,\%$ at every output time:
$\langle\vv,\mathbf E\rangle$ is negligible. The mass change therefore enters
through something that modifies $\xi$ \emph{outside} its prognostic equation.

\section{The leak is the $l_q$ smoothing of the nonlinear $\xi$}

The only other operation that touches $\xi$ is the transform. With $kDim=99$
collocation points and $b_{kDim}=66$ cubic B--spline coefficients
($\mathrm{mubar}=3$ Gauss nodes per cell), \verb|spectralTransform!| applies, on
every step, a regularized least--squares fit $P_{l_q}$ that minimizes
\begin{equation}
  \big\|\,f-\text{data}\,\big\|^2_{W} \;+\; l_q \int \big(f'''\big)^2 ,
  \label{eq:lq}
\end{equation}
where $W$ are the Gauss weights and the second term (the Springsteel $Q$--matrix,
$l_q=2$ by default) suppresses high--wavenumber overshoot for stability. The model
then re--evaluates the field from this fit each step
($\xi^n\to P_{l_q}\xi^{\star}$).

\paragraph{The bare projection conserves; the smoothing of a nonlinear variable
does not.} With $l_q=0$, \eqref{eq:lq} is the plain Gauss--weighted Galerkin
projection, which preserves the integral of the projected field (constants lie in
the spline span, and the $\mathrm{mubar}=3$ nodes integrate cubics exactly). So
$b_{kDim}<kDim$ alone is \emph{not} the problem. The trouble is that the model
projects the \emph{nonlinear} $\xi=\ln(\rd/\rho_0)$ \emph{with} $l_q>0$. The
smoothing lowers the variance of $\xi$ while approximately preserving $\int\xi$;
because $e^{(\cdot)}$ is convex, Jensen's inequality gives
\begin{equation}
  \int \rho_0\,e^{\,P_{l_q}\xi} \;<\; \int \rho_0\,e^{\xi},
  \label{eq:jensen}
\end{equation}
a \emph{systematic} (always--negative) loss of $\int\rd$. Projecting the
\emph{linear} density $\rd$ instead has no such convexity bias: $P_{l_q}$
preserves $\int\rd$ directly.

\paragraph{Direct confirmation.} Iterating the actual update with the velocity
frozen at $t=200$\,s --- so the conservative transport is held fixed and only the
fit acts --- isolates the projection. Two knobs (project $\xi$ vs.\ $\rd$; $l_q=2$
vs.\ $0$) give:
\begin{center}
\begin{tabular}{llrr}
\hline
variable projected & filter & loss rate [kg/s] & fraction of observed \\
\hline
$\xi$ (log--density)   & $l_q=2.0$ & $-21.8$            & $0.99$ \\
$\xi$ (log--density)   & $l_q=0.0$ & $-1.8\times10^{-1}$ & $0.008$ \\
$\rd$ (linear / flux)  & $l_q=2.0$ & $\phantom{-}4\times10^{-7}$ & $\approx 0$ \\
$\rd$ (linear / flux)  & $l_q=0.0$ & $\phantom{-}3\times10^{-7}$ & $\approx 0$ \\
\hline
\end{tabular}
\end{center}
The $\xi$/$l_q{=}2$ loop reproduces $99\,\%$ of the observed model rate
($-22.1$\,kg/s). Turning the smoothing off ($l_q=0$) cuts the loss $120\times$.
Projecting the linear density conserves to $\sim10^{-15}$ regardless of $l_q$.
For contrast, the dry run's frozen loop loses $\sim10^{-2}$\,kg/s,
$\sim2000\times$ less, because its smooth $\xi'$ carries almost no high--wavenumber
content for the penalty to act on. A single offline application of $P_{l_q}$ to one
instantaneous tendency captures only $\sim7\,\%$ of the drift: the loss is a
\emph{steady--state} balance in which the RHS continually regenerates
high--wavenumber $\xi$ structure and the smoothing continually removes it.

\section{Why the moist run, and the link to the reference imbalance}

The high--wavenumber $\xi$ content that the penalty acts on is generated mainly by
$-w\,\partial\xhat/\partial z$ in \eqref{eq:cont}: vertical motion sliding $\xi$
along the steep background density gradient $\partial\xhat/\partial z$ (the
largest gradient in the problem). The rate of structure generation --- and hence
the smoothing loss --- therefore scales with the velocity field. In the dry case
the flow is a weak, transient, closed overturning and $\xi'$ stays smooth. In the
moist case the unbalanced base--state cloud loading (the reference density
$\bar\rho$ omits the condensate weight that $\rho_t$ carries via $\mu_c$) drives a
strong, persistent, nearly domain--wide circulation; that flow continually
injects high--wavenumber structure into $\xi$, which the $l_q$ term then damps via
\eqref{eq:jensen} every step. The Chebyshev (Z--grid) variant, which has no such
spline smoothing, drifts only $\sim10^{-4}\,\%$, corroborating a basis$+$filter
cause.

This reconciles the two layers: the reference imbalance does not lose mass by
itself --- it sets up the velocity, and the velocity feeds the smoothing loss. A
balanced initial state (condensate weight included in $\bar\rho$) removes the
circulation and so removes the loss, which is why it passes the benchmark; but it
is a trigger fix, not a cure for the underlying behaviour.

\section{Implications}

\begin{itemize}
\item The non--conservation is the $l_q$ third--derivative smoothing acting on the
\emph{nonlinear} log--density $\xi$ (Jensen bias, \eqref{eq:jensen}), not the
advective form of continuity, not $b_{kDim}<kDim$ by itself, and not a
boundary--flux error.
\item \textbf{Structural cure:} carry/project the \emph{linear} density $\rd$
(flux form). The smoothing then preserves $\int\rd$ directly --- confirmed to
$\sim10^{-15}$ \emph{with} $l_q=2$ --- so stability and conservation coexist.
\item Setting $l_q=0$ removes $\sim99\,\%$ of the loss but invites the spectral
overshoot the penalty exists to suppress, so it is not a usable fix on its own.
\item For idealized balanced benchmarks, removing the spurious circulation (e.g.
including the base condensate loading in $\bar\rho$) suppresses the trigger and
recovers near machine--precision conservation.
\end{itemize}

The offline diagnostic that produced every number here is
\verb|benchmarks/mass_error_diagnostic.jl| (\verb|moist|/\verb|dry|); it changes
no model code and uses the model's own spectral operators.

\end{document}
